

\chapter{Exercicis proposats}
\label{cap:exercicis-proposats}

\section{Els nombres naturals}

\begin{exercici}[Ús dels axiomes de Peano]
Usant només els axiomes de Peano, prova que
\[
2\neq 0,\qquad 2\neq 1,\qquad 5\neq 3.
\]
\end{exercici}

\begin{indicacio}
Escriu \(1=S(0)\), \(2=S(1)\), \(3=S(2)\), etc. Usa que \(0\) no és successor de cap natural i que el successor és injectiu.
\end{indicacio}


\begin{exercici}[Suma dels primers naturals]
Demostra per inducció que, per a tot \(n\in\N\),
\[
2(0+1+2+\cdots+n)=n(n+1).
\]
\end{exercici}

\begin{indicacio}
Fes inducció sobre \(n\). En el pas inductiu, escriu la suma fins a \(n+1\) com la suma fins a \(n\) més \(n+1\), i factoritza \(n+1\).
\end{indicacio}


\begin{exercici}[Suma dels quadrats]
Demostra per inducció que, per a tot \(n\in\N\), amb \(n\geq 1\),
\[
6(1^2+2^2+\cdots+n^2)=n(n+1)(2n+1).
\]
\end{exercici}

\begin{indicacio}
Fes inducció sobre \(n\). En el pas inductiu, afegeix \(6(n+1)^2\) als dos membres i comprova que el resultat es pot factoritzar com \((n+1)(n+2)(2n+3)\).
\end{indicacio}


\begin{exercici}[Cancel\textperiodcentered{}lació de la suma]
Demostra que si \(a,b,c\in\N\) i
\[
a+c=b+c,
\]
aleshores
\[
a=b.
\]
\end{exercici}

\begin{indicacio}
Fes inducció sobre \(c\). En el pas inductiu, transforma \(a+S(c)=b+S(c)\) en \(S(a+c)=S(b+c)\) i usa la injectivitat del successor.
\end{indicacio}


\begin{exercici}[Cancel\textperiodcentered{}lació de l'ordre amb la suma]
Demostra que, per a tots \(a,b,t\in\N\), si
\[
a+t\leq b+t,
\]
aleshores
\[
a\leq b.
\]
\end{exercici}

\begin{indicacio}
Usa la definició d'ordre: de \(a+t\leq b+t\) obtens \(b+t=(a+t)+n\) per a algun \(n\in\N\). Reordena els termes i aplica la cancel\textperiodcentered{}lació de la suma.
\end{indicacio}


\begin{exercici}[Compatibilitat de l'ordre amb la suma]
Demostra que si \(a,b,c\in\N\) i
\[
a\leq b,
\]
aleshores
\[
a+c\leq b+c.
\]
\end{exercici}

\begin{indicacio}
Escriu \(b=a+n\) per a algun \(n\in\N\). Després suma \(c\) i reordena:
\[
b+c=(a+c)+n.
\]
\end{indicacio}


\begin{exercici}[Compatibilitat de l'ordre amb el producte]
Demostra que si \(a,b,c\in\N\) i
\[
a\leq b,
\]
aleshores
\[
a\cdot c\leq b\cdot c.
\]
A més, si \(a<b\) i \(c\neq0\), demostra que
\[
a\cdot c<b\cdot c.
\]
\end{exercici}

\begin{indicacio}
Escriu \(b=a+n\). En multiplicar per \(c\), obtens
\[
b\cdot c=a\cdot c+n\cdot c.
\]
Per a la desigualtat estricta, usa que si \(n\neq0\) i \(c\neq0\), aleshores \(n\cdot c\neq0\).
\end{indicacio}


\begin{exercici}[Cancel\textperiodcentered{}lació del producte]
Demostra que si \(a,b,c\in\N\), \(c\neq0\), i
\[
a\cdot c=b\cdot c,
\]
aleshores
\[
a=b.
\]
\end{exercici}

\begin{indicacio}
Usa la tricotomia de l'ordre. Si \(a<b\), aleshores \(a\cdot c<b\cdot c\); si \(b<a\), aleshores \(b\cdot c<a\cdot c\). Tots dos casos contradiuen la igualtat.
\end{indicacio}


\begin{exercici}[Potències]
Definim les potències d'un natural \(a\) per recursió:
\[
a^0=1,
\qquad
a^{S(n)}=a^n\cdot a.
\]
Demostra que, per a tots \(m,n\in\N\),
\[
a^m\cdot a^n=a^{m+n}.
\]
\end{exercici}

\begin{indicacio}
Fixa \(m\) i fes inducció sobre \(n\). En el pas inductiu, usa
\[
a^{S(n)}=a^n\cdot a
\]
i l'associativitat del producte.
\end{indicacio}


\begin{exercici}[Algorisme d'Euclides]
Calcula
\[
\mcd(924,396)
\]
amb l'algorisme d'Euclides i expressa'l com a combinació lineal amb coeficients enters de \(924\) i \(396\).
\end{exercici}

\begin{indicacio}
Comença dividint \(924\) entre \(396\):
\[
924=396\cdot2+132.
\]
Després divideix \(396\) entre \(132\). Per a la combinació lineal, aïlla el residu \(132\) de la primera igualtat.
\end{indicacio}


\begin{exercici}[Descomposició en factors primers]
Descompon en factors primers els nombres \(840\), \(1155\) i \(2025\). Calcula després el \(\mcd\) i el \(\mcm\) de \(840\) i \(1155\).
\end{exercici}

\begin{indicacio}
Descompon cada nombre en potències de primers. Per al \(\mcd\), pren els factors comuns amb l'exponent més petit. Per al \(\mcm\), pren tots els factors amb l'exponent més gran.
\end{indicacio}


\section{Els nombres enters}

\begin{exercici}[Classes d'equivalència]
Comprova si els parells següents són equivalents en la construcció de \(\Z\):
\[
(8,13),\qquad (2,7),\qquad (10,14),\qquad (15,20).
\]
Agrupa'ls per classes d'equivalència.
\end{exercici}

\begin{indicacio}
Recorda que
\[
(a,b)\sim(c,d)\iff a+d=b+c.
\]
Compara cada parell amb \((8,13)\). Tres dels parells representen el mateix enter.
\end{indicacio}


\begin{exercici}[Representants d'una mateixa classe]
Demostra que, per a tots \(a,b,k\in\N\),
\[
\classe{(a,b)}=\classe{(a+k,b+k)}.
\]
\end{exercici}

\begin{indicacio}
Cal provar que
\[
(a,b)\sim(a+k,b+k).
\]
És a dir, comprova que
\[
a+(b+k)=b+(a+k)
\]
usant l'associativitat i la commutativitat de la suma en \(\N\).
\end{indicacio}


\begin{exercici}[Suma, resta i producte]
Calcula, com a classes d'equivalència:
\[
\classe{(6,11)}+\classe{(9,2)},
\qquad
\classe{(6,11)}-\classe{(9,2)},
\qquad
\classe{(6,11)}\cdot\classe{(9,2)}.
\]
Interpreta cada resultat en la notació usual dels enters.
\end{exercici}

\begin{indicacio}
Per a la suma, suma component a component. Per a la resta, primer pren l'oposat de \(\classe{(9,2)}\). Per al producte, usa
\[
\classe{(a,b)}\cdot\classe{(c,d)}
=
\classe{(ac+bd,ad+bc)}.
\]
\end{indicacio}


\begin{exercici}[L'oposat d'un enter]
Prova que l'oposat de
\[
\classe{(a,b)}
\]
és
\[
\classe{(b,a)}.
\]
\end{exercici}

\begin{indicacio}
Suma les dues classes:
\[
\classe{(a,b)}+\classe{(b,a)}.
\]
Has d'obtenir una classe equivalent a \(\classe{(0,0)}\).
\end{indicacio}


\begin{exercici}[Equacions de la forma \(x+a=b\)]
Resol en \(\Z\) l'equació
\[
x+7=5.
\]
Dona la solució com a classe d'equivalència i en la notació usual dels enters.
\end{exercici}

\begin{indicacio}
Aïlla formalment:
\[
x=5-7.
\]
Després escriu \(5=\classe{(5,0)}\), \(7=\classe{(7,0)}\) i usa la definició de resta.
\end{indicacio}


\begin{exercici}[Producte d'enters i signe]
Calcula els productes següents i interpreta'n el signe:
\[
\classe{(0,3)}\cdot\classe{(0,4)},
\qquad
\classe{(5,0)}\cdot\classe{(0,2)},
\qquad
\classe{(0,6)}\cdot\classe{(7,0)}.
\]
\end{exercici}

\begin{indicacio}
Interpreta primer les classes com enters usuals:
\[
\classe{(0,3)}=-3,\quad \classe{(0,4)}=-4,\quad \classe{(5,0)}=5.
\]
Després comprova els resultats amb la fórmula del producte de classes.
\end{indicacio}


\begin{exercici}[Ordre en \(\Z\)]
Usant la definició d'ordre en \(\Z\), ordena de menor a major les classes següents, indicant si n'hi ha dues d'iguals:
\[
\classe{(1,8)},
\qquad
\classe{(5,5)},
\qquad
\classe{(9,2)},
\qquad
\classe{(3,10)}.
\]
\end{exercici}

\begin{indicacio}
Pots interpretar cada classe com la diferència entre els seus components. Observa especialment que \(\classe{(1,8)}\) i \(\classe{(3,10)}\) representen el mateix enter.
\end{indicacio}


\begin{exercici}[Compatibilitat de l'ordre amb la suma]
Demostra que, si \(x\leq y\) en \(\Z\), aleshores
\[
x+z\leq y+z
\]
per a tot \(z\in\Z\).
\end{exercici}

\begin{indicacio}
Usa la caracterització
\[
x\leq y\iff y-x\geq0.
\]
Després observa que
\[
(y+z)-(x+z)=y-x.
\]
\end{indicacio}


\begin{exercici}[Multiplicació d'una desigualtat]
Siguin \(a,b,c\in\Z\), amb
\[
a\leq b.
\]
Demostra que:
\[
\begin{array}{ll}
\text{\textbf{(i)}} & \text{si } c\geq0,\text{ aleshores } ac\leq bc,\\[1mm]
\text{\textbf{(ii)}} & \text{si } c\leq0,\text{ aleshores } bc\leq ac.
\end{array}
\]
\end{exercici}

\begin{indicacio}
De \(a\leq b\) dedueix \(b-a\geq0\). Si \(c\geq0\), estudia \((b-a)c\). Si \(c\leq0\), estudia \((b-a)(-c)\).
\end{indicacio}


\begin{exercici}[Resolució d'una desigualtat]
Resol en \(\Z\) la desigualtat
\[
-3x+2\leq 11.
\]
\end{exercici}

\begin{indicacio}
Resta \(2\) als dos membres. Després multiplica per \(-1\) i recorda que el sentit de la desigualtat s'inverteix.
\end{indicacio}


\begin{exercici}[Càlculs amb valor absolut]
Calcula:
\[
|-8|,
\qquad
|5-12|,
\qquad
|-4\cdot6|.
\]
\end{exercici}

\begin{indicacio}
Calcula primer l'enter que hi ha dins del valor absolut i després pren-ne la distància a \(0\).
\end{indicacio}


\begin{exercici}[Oposat i valor absolut]
Demostra que, per a tot \(n\in\Z\),
\[
|-n|=|n|.
\]
\end{exercici}

\begin{indicacio}
Distingueix els casos \(n\geq0\) i \(n<0\). En cada cas, aplica directament la definició de valor absolut.
\end{indicacio}


\begin{exercici}[Distància entre enters]
Demostra que, per a tots \(n,m\in\Z\),
\[
|n-m|=|m-n|.
\]
\end{exercici}

\begin{indicacio}
Observa que
\[
n-m=-(m-n).
\]
Després aplica la propietat \(|-x|=|x|\).
\end{indicacio}


\begin{exercici}[Valor absolut d'un producte]
Demostra que, per a tots \(n,m\in\Z\),
\[
|nm|=|n||m|.
\]
\end{exercici}

\begin{indicacio}
Distingueix quatre casos segons els signes de \(n\) i \(m\): tots dos no negatius, un no negatiu i l'altre negatiu, o tots dos negatius.
\end{indicacio}


\begin{exercici}[Desigualtat triangular]
Demostra que, per a tots \(n,m\in\Z\),
\[
|n+m|\leq |n|+|m|.
\]
\end{exercici}

\begin{indicacio}
Parteix de
\[
-|n|\leq n\leq |n|
\qquad\text{i}\qquad
-|m|\leq m\leq |m|.
\]
Suma les dues desigualtats i aplica la caracterització de \(|x|\leq k\).
\end{indicacio}


\begin{exercici}[Desigualtats amb valor absolut]
Troba tots els enters \(n\in\Z\) que compleixen
\[
|n-3|\leq2.
\]
\end{exercici}

\begin{indicacio}
Transforma la desigualtat en
\[
-2\leq n-3\leq2.
\]
Després suma \(3\) als tres membres.
\end{indicacio}


\begin{exercici}[Teorema de la divisió en \(\Z\)]
Troba el quocient i el residu de la divisió de \(-37\) entre \(5\), i també de la divisió de \(-37\) entre \(-5\).
\end{exercici}

\begin{indicacio}
El residu ha de ser sempre no negatiu i menor que el valor absolut del divisor. Busca \(q\) de manera que
\[
-37=bq+r,\qquad 0\leq r<|b|.
\]
\end{indicacio}


\begin{exercici}[Divisibilitat en \(\Z\)]
Determina si les afirmacions següents són certes o falses:
\[
-4\mid20,
\qquad
4\mid(-20),
\qquad
-5\mid(-35),
\qquad
6\mid20.
\]
\end{exercici}

\begin{indicacio}
Usa la definició: \(a\mid b\) si existeix \(q\in\Z\) tal que \(b=aq\). En cada cas, intenta trobar aquest enter \(q\).
\end{indicacio}


\begin{exercici}[Divisors enters]
Troba tots els divisors enters de \(18\).
\end{exercici}

\begin{indicacio}
Comença trobant els divisors positius de \(18\). Després afegeix els oposats corresponents, perquè en \(\Z\) els divisors apareixen per parelles de signe.
\end{indicacio}


\begin{exercici}[Diferència de múltiples]
Siguin \(a,b,c\in\Z\), amb \(a\neq0\). Demostra que, si
\[
a\mid b
\qquad\text{i}\qquad
a\mid c,
\]
aleshores
\[
a\mid(b-c).
\]
\end{exercici}

\begin{indicacio}
Escriu \(b=aq\) i \(c=ar\), amb \(q,r\in\Z\). Després factoritza:
\[
b-c=aq-ar=a(q-r).
\]
\end{indicacio}


\begin{exercici}[Transitivitat de la divisibilitat]
Siguin \(a,b,c\in\Z\), amb \(a\neq0\) i \(b\neq0\). Demostra que, si
\[
a\mid b
\qquad\text{i}\qquad
b\mid c,
\]
aleshores
\[
a\mid c.
\]
\end{exercici}

\begin{indicacio}
Escriu \(b=aq\) i \(c=br\), amb \(q,r\in\Z\). Substitueix \(b\) dins de l'expressió de \(c\).
\end{indicacio}

